\begin{corollary}\label{cor:semipositive-dual-lp-limit}
Under the hypotheses of \cref{thm:semipositive-dual-exact-equation}, suppose in addition that $f\in L^p(\mathrm{d}V)$ for some $p>1$. Then the normalized threshold solution $u_0$ is continuous, and the normalized greatest exponential solutions converge uniformly to $u_0$ as the exponent tends to zero.
\end{corollary}
\begin{proof}
By \cref{prop:individual-normalizer-lp}, the constant $c$ for which $\mathrm{MA}_\omega(v)=cf\,\mathrm{d}V$ has a continuous solution is $b_\omega(f)$. The continuous normalized solution is therefore a threshold solution, and it equals $u_0$ by the uniqueness in \cref{thm:semipositive-dual-exact-equation}. For $L^p$ data the greatest exponential solution coincides with the bounded solution, as explained after \cref{thm:individual-exponential-subsolution}. By \cref{prop:symmetric-equation-clusters}, applied with the trivial group, the normalized family is uniformly precompact, and every cluster limit solves the threshold equation. Uniqueness identifies every cluster limit with $u_0$, which proves uniform convergence of the whole family.
\end{proof}
